\chapter*{INTRODUCTION}
\addcontentsline{toc}{chapter}{INTRODUCTION}

\section*{Rationale}

Vietnamese online news has become a primary medium of public information, with
tens of millions of articles published each year across major outlets such as
\textit{tuoitre.vn}, \textit{thanhnien.vn}, \textit{tienphong.vn},
\textit{vietnamnet.vn}, etc. This volume creates two parallel
needs for the readers. The first is \textbf{compression}: a quick,
trustworthy summary that lets a reader skim a long article in a few seconds.
The second is \textbf{verification}: a check on whether the article's main
claims are supported by other reputable sources, particularly as
misinformation, clickbait, and partial reporting circulate alongside accurate
journalism.

Currently, Vietnamese summarizers built on a single Large Language Model
(LLM) still produce uneven outputs --- sometimes too long, sometimes too
literal, occasionally hallucinated. A recent line of work,
\textbf{Mixture-of-Agents} (MoA) by Wang et al.\ \cite{Wang2024}, argues that
multiple LLMs can be combined in a hierarchical proposer--aggregator structure
to outperform any single model on instruction-following benchmarks. Whether
this benefit transfers to \textit{grounded} Vietnamese summarization --- where
the system must remain faithful to a source article rather than answer
free-form prompts --- is an open question that has not been answered in the
Vietnamese NLP literature.

A second, deeper problem motivates this thesis. The metrics most commonly used
to evaluate summarization in Vietnamese research --- ROUGE \cite{Lin2004},
BLEU, and BERTScore \cite{Zhang2020} computed against the source article ---
measure overlap, not quality. They reward summaries that reuse phrases from
the source and punish summaries that paraphrase or compress aggressively.
A correctly-implemented aggregator that synthesizes drafts into a cleaner
editorial summary will, by construction, score lower on these metrics than a
draft that copies surface forms. If Vietnamese summarization research is to
adopt aggregator-style systems, it needs an evaluation methodology that does
not structurally disadvantage them.

\section*{Scientific and Practical Significance}

\textbf{Scientific significance.}
This thesis contributes a \textit{three-axis evaluation framework} that
combines (i)~content-retention metrics against the source, (ii)~LLM-as-judge
quality and pairwise preference scoring, and (iii)~blind human ranking with
inter-rater agreement statistics. It uses this framework to expose, with
empirical evidence on Vietnamese data, a structural limitation of overlap
metrics for aggregator-style summarization. This limitation has been mentioned indirectly
to in English-language work \cite{Wang2024,Zheng2023} but has not been
explicitly studied for Vietnamese.

\textbf{Practical significance.}
The thesis ships an end-to-end product, called \textbf{Fiber}, that any
Vietnamese reader can install in a Chromium-based browser. Fiber injects a
sidebar into seven major Vietnamese news domains, summarizes the article in
real time using a chosen LLM (or a fused MoA pipeline), and fact-checks the
summary against live web evidence via the Tavily search API. The whole
system --- a Next.js backend, a Plasmo browser extension, and a Python
BERTScore microservice --- is open-source and reproducible.

\section*{Research Objectives}

\begin{enumerate}
  \item Design and build an end-to-end browser-extension system that
        summarizes and fact-checks Vietnamese news articles in real time,
        supporting multiple LLM providers (OpenAI, Anthropic, Google) and
        multiple routing strategies.
  \item Adapt the Mixture-of-Agents pipeline of Wang et al.\ (2024) to the
        grounded Vietnamese summarization setting, including a source-article
        residual connection in the aggregator prompt.
  \item Propose and validate a three-axis evaluation framework --- content
        retention (Axis~A), LLM-judge quality and preference (Axis~B), and
        human ranking with agreement statistics (Axis~C) --- and use it to
        compare MoA fusion against single-model baselines on a topic-balanced
        Vietnamese dataset.
\end{enumerate}

\section*{Scope and Object of Study}

\textbf{Object.} Vietnamese-language news articles published on
\textit{tuoitre.vn}, \textit{thanhnien.vn}, \textit{tienphong.vn},
\textit{vietnamnet.vn}, \textit{laodong.vn}, \textit{vtv.vn}, and
\textit{nld.com.vn}. The empirical evaluation in Chapter~4 uses 154 articles
from \textit{tienphong.vn}, balanced across five topics: current affairs
(\textit{thời sự}), law (\textit{pháp luật}), economy (\textit{kinh tế}),
education (\textit{giáo dục}), and culture (\textit{văn hóa}).

\textbf{Scope.}
\begin{itemize}
  \item \textit{Modeling.} Three commercial cloud LLMs (GPT-4o-mini, Claude
        Haiku 4.5, Gemini 2.5 Flash) as proposers; GPT-4o as the aggregator;
        best draft selection as the strongest single-model baseline.
        Vietnamese-specific models (ViT5 \cite{Nguyen2022} and
        PhoBERT \cite{Nguyen2020}) are used for
        routing classification and embedding-based scoring.
  \item \textit{Evaluation.} Single domain (\textit{tienphong.vn}), single run
        per article, no multi-domain generalization claim. The human study is
        comprehensive with $n = 17$ raters and $48$ tasks, framed as a descriptive
        complement to Axis~A and Axis~B.
  \item \textit{Fact-checking.} Search-augmented verification via the Tavily
        web-search API. Claim-detection and evidence-retrieval are
        implemented; deeper structured-knowledge fact-checking
        (knowledge-graph or stance-classification) is out of scope.
\end{itemize}

\section*{Research Methodology}

The thesis follows a comparative experimental methodology. The same $154$
articles are summarized by five approaches (three individual proposers, the
MoA fused output, and the best draft single-model baseline), and the resulting
$770$ summaries are evaluated along three axes. Axis~A uses ROUGE-1/2/L,
BLEU, BERTScore (with PhoBERT) and compression rate; Axis~B uses a
FLASK-derived 5-dimension rubric \cite{Ye2023}, an MT-Bench-style absolute
score \cite{Zheng2023}, an AlpacaEval-style pairwise preference judge with
position randomization \cite{Dubois2024}, and a claim-level factuality
classifier; Axis~C uses a blind $K = 3$ ranking interface administered
through the same browser-extension product. Statistical tests include the
two-sided sign test (ties excluded) for pairwise win rates and Fleiss'
$\kappa$ \cite{Landis1977} for inter-rater agreement; length bias is
controlled via the simplified-bucket method of
Dubois et al.\ \cite{Dubois2024}.

\section*{Outline of the Thesis}

The remainder of this thesis is organized as follows.

\textbf{Chapter 1 --- Background.} Introduces the summarization and fact-checking
problems, the Mixture-of-Agents technique, and related Vietnamese and
international work, identifying the gap that this thesis fills.

\textbf{Chapter 2 --- Theoretical Foundations.} Reviews the LLM architectures used
in the system, the three families of evaluation metrics (overlap, LLM-judge,
human-ranking), the factuality methodology, and the technologies used to
build the browser extension and back-end services.

\textbf{Chapter 3 --- System Design and Implementation.} Presents requirement
analysis, the three-service architecture, the multi-provider LLM abstraction,
the MoA pipeline (the technical contribution), the three-axis evaluation
framework (the methodological contribution), the database schema, and the
user interface.

\textbf{Chapter 4 --- Experiments and Evaluation.} Reports results along the three
axes on the 154-article dataset, including the headline pairwise result
(fused beats best draft $68.9\%$, $p < 0.0001$), the topic-balanced
breakdown, the human study interpretation, the cross-axis triangulation, and
the methodological limitations.

The \textbf{Conclusion} summarizes contributions, limitations, and future work;
three appendices document examples, the aggregator-prompt ablation, and
the fact-checking pipeline.
